% !TEX root = ../../../proposal.tex

\subsection{Compositional Diffusion Models}
\label{sec:compositional_diffusion}

Single text-to-image diffusion models often fail when a prompt requires several factors to
be combined correctly. Recent analyses \citep{t2i_compbench_2025, du2024compositional,
mechanisms2025projective} show consistent errors in attribute binding, spatial relations,
and numeracy, even when the image remains visually plausible. These failures arise because
the models largely learn correlations in joint image-text data rather than an explicit
factorized structure over objects, attributes, and relations \citep{t2i_compbench_2025}.
Compositional generative modeling aims to overcome this limitation by generating new factor
combinations in a structured and semantically meaningful way. Earlier work framed this in
terms of correctness, coverage, and compositionality \citep{vedantam2017generative}, and
models such as SCAN \citep{higgins2017scan} showed the value of controllable latent
structure. However, the field still lacks a clear account of what a composition should
preserve, why some definitions fail, and what structural conditions make composition
possible \citep{mechanisms2025projective, du2024compositional}.


The next three definitions review the main formulations of composition in this literature: the \emph{simple product}, \emph{Bayes composition}, and \emph{projective composition}. 
The first two define composition by combining densities directly, while the third defines it through what the composed result should preserve.

\subsubsection{Simple Product}

The first definition is the \emph{simple product}. Given two distributions \(p_1\) and
\(p_2\) over \(\mathbb{R}^n\), it defines the composition as
\[
    \hat p(x) \propto p_1(x)p_2(x).
\]
This is a familiar way to combine distributions, but it has an important limitation: the
composed distribution cannot be truly out-of-distribution with respect to \(p_1\) or
\(p_2\). If either \(p_1(x)=0\) or \(p_2(x)=0\), then \(\hat p(x)=0\) as well. As
\citet{mechanisms2025projective} show, this becomes a problem when the desired
composition lies outside the support of both source distributions. In their CLEVR example,
if each source distribution contains only one object, then any image with \emph{both}
objects has zero probability under both \(p_1\) and \(p_2\), so the simple product also
assigns it zero probability. 

\subsubsection{Bayes Composition}

The second prior definition is \emph{Bayes composition}. Here the idea is not to multiply
arbitrary densities, but to interpret the source models as conditional distributions under a
shared data distribution. If \(c_1\) and \(c_2\) denote two conditions of interest, then
the desired composition is defined as the conditional distribution
\[
    \hat p(x) := p(x \mid c_1, c_2).
\]
Assuming \(c_1\) and \(c_2\) are conditionally independent given \(x\), this gives
\begin{equation}
    \hat p(x)
    = p(x \mid c_1, c_2)
    \propto \frac{p(x \mid c_1)\,p(x \mid c_2)}{p(x)}.
    \label{eq:poe_dist}
\end{equation}
This is the Bayes composition used in compositional diffusion work such as
\citep{liu2022compositional,du2023reduce,du2024compositional,NEURIPS2020_49856ed4}.
It works best when the desired composition already corresponds to a valid conditional distribution under the
training data, that is, when the component concepts have been seen together often enough
for the model to treat their conjunction as part of the same underlying distribution. In
that setting, each conditional contributes evidence for the target sample, while the
unconditional term removes shared evidence that would otherwise be counted twice. This is
the setting that motivates the co-occurring concept pairs in Group~1 of Phase~1.

However, Bayes composition is not truly out-of-distribution \citep{mechanisms2025projective}. 
If the desired conjunction is not supported by the underlying data distribution, then the composition cannot create it
simply by combining conditionals. It therefore succeeds when the
joint concept is already meaningful under the training distribution, and fails when the
desired combination lies outside that support.

To use Bayes composition in diffusion models, we convert the composed distribution into a
composed score. Starting from
\[
    \hat p(x \mid c_1, c_2)
    \propto
    \frac{p(x \mid c_1)\,p(x \mid c_2)}{p(x)},
\]
take the gradient of the log on both sides. This gives
\begin{equation}
    \nabla_x \log \hat p(x \mid c_1, c_2)
    =
    \nabla_x \log p(x \mid c_1)
    +
    \nabla_x \log p(x \mid c_2)
    -
    \nabla_x \log p(x).
    \label{eq:poe_score}
\end{equation}
The composed score is therefore the sum of the two conditional scores minus the
unconditional score.

Earlier, the DDPM score was written in terms of the noise predictor as
\[
    s_\theta(x_t,t,\cdot)
    =
    -\frac{\varepsilon_\theta(x_t,t,\cdot)}{\sqrt{1-\bar\alpha_t}}.
\]
This means that the same composition can be implemented directly in terms of noise
predictions. Let \(\varepsilon_\theta(x_t,t,\varnothing)\) denote the unconditional
prediction and \(\varepsilon_\theta(x_t,t,c_i)\) the prediction conditioned on concept
\(c_i\). For a single condition, classifier-free guidance strengthens the effect of that
condition relative to the unconditional background through
\[
    \varepsilon_\theta(x_t,t,\varnothing)
    +
    w_i\bigl[\varepsilon_\theta(x_t,t,c_i)-\varepsilon_\theta(x_t,t,\varnothing)\bigr].
\]
Applying this idea to both concepts gives the practical AND operator used by
\citet{liu2022compositional}:
\begin{equation}
    \tilde\varepsilon_{\mathrm{AND}}
    =
    \varepsilon_\theta(x_t,t,\varnothing)
    +
    w_1\bigl[\varepsilon_\theta(x_t,t,c_1)-\varepsilon_\theta(x_t,t,\varnothing)\bigr]
    +
    w_2\bigl[\varepsilon_\theta(x_t,t,c_2)-\varepsilon_\theta(x_t,t,\varnothing)\bigr].
    \label{eq:and_operator}
\end{equation}
When \(w_1=w_2=1\), this reduces to
\[
    \varepsilon_\theta(x_t,t,c_1)
    +
    \varepsilon_\theta(x_t,t,c_2)
    -
    \varepsilon_\theta(x_t,t,\varnothing),
\]
which is exactly the Bayes composition written in noise-prediction form. This is the
practical score-space composition rule used later in Phase~1 and Phase~3.



\subsubsection{Projective Composition}
\label{subsec:bg_projective}

\citet{mechanisms2025projective} argue that a useful definition of composition must answer
two separate questions: what the composed distribution should preserve, and when a practical
compositional method can realize it. Their notion of \emph{projective composition} addresses
the first question. Under this view, the composed distribution is judged by the part each
source is meant to contribute. For example, if one source determines style and another
content, then a composition is correct if it preserves the style of the first and the
content of the second, even when neither source distribution contains that full combination.
This is closely aligned with \citet{du2024compositional}, where composition is understood as
building a new distribution from simpler components that each contribute part of the final
structure.

This matters because it allows composition that is out of distribution with respect to the
source distributions being combined. The final sample can combine properties from different
sources without needing to belong to any one source distribution as a whole. 
As Figure~\ref{fig:projective_composition} illustrates, the composed distribution can differ
from both source distributions in the original sample space while still matching each source
after applying the corresponding projection.
That is precisely what the simple product cannot express, since it is confined to overlapping
support, and what Bayes composition cannot justify when the desired conjunction is not
already supported by some underlying joint distribution.

\begin{figure}[H]
    \centering
    \includegraphics[width=0.58\linewidth]{chapters/background/media/projective_composition.png}
    \caption{Illustration of projective composition. The composed distribution \(\hat p\) need not coincide with either source distribution \(p_1\) or \(p_2\) in the original sample space. Instead, correctness is defined after applying source-specific projections: under \(\Pi_1\), the composed distribution matches the aspect preserved from \(p_1\), and under \(\Pi_2\), it matches the aspect preserved from \(p_2\). The figure therefore shows how a composition can be out of distribution with respect to the sources as whole distributions while still preserving the designated contribution of each source.}
    \label{fig:projective_composition}
\end{figure}

The key implication is representational. 
Standard disentanglement is not enough if concepts still interfere when one is changed. What composition requires is a stronger separation in which each concept acts on its own factor without perturbing the others. 
\citet{mechanisms2025projective} further show that this separation need not appear in pixel space itself. 
If it exists in some invertible feature space, then the same composition operator yields a correct projective composition with respect to those feature-space factors. 
This is a statement about when the target composition is well defined, not a guarantee that reverse-diffusion sampling will always recover it. 
For the proposed study, the lesson is that compositional success depends on whether the learned representation allows the concepts to be preserved independently.
